\documentclass[14pt]{extarticle}
\usepackage{knu-fcc-style}
\usepackage{booktabs}

\begin{document}

\KNUTitlePageMake[
  type={Звіт до лабораторної роботи №4},
  specialty={122 Комп'ютерні науки},
  title={Обробка відео на GPU із CUDA},
  studentInfo={Виконав студент}, student={Божевський Іван},
  studentSignature={false}, supervisorSignature={false}, year={2026}
]

\KNUUnnumberedSection{Реферат}
У роботі реалізовано конвеєр обробки відео засобами OpenCV і CUDA. Відеофайл декодується послідовно, поточний та попередній кадри передаються до пам'яті GPU, для них обчислюється абсолютна попарна різниця, а результат записується до нового MP4-файлу. Програма не має CPU fallback і потребує OpenCV, скомпільованого з підтримкою CUDA.

\KNUKeywords{CUDA, GPU, OPENCV, ОБРОБКА ВІДЕО, ПОПАРНА РІЗНИЦЯ КАДРІВ}
\KNUContents

\KNUUnnumberedSection{Вступ}
Метою роботи є реалізація обробки відео на графічному процесорі. Завдання включає відкриття відеофайлу за допомогою OpenCV, завантаження кадрів у пам'ять GPU, обчислення попарної різниці кадрів та збереження результуючого відео засобами OpenCV. Така операція є базою для виявлення руху, фоновіднімання та попередньої обробки відеопотоку.

\section{Теоретичні відомості та алгоритм}
Нехай $I_t(x,y,c)$ --- значення каналу $c$ пікселя з координатами $(x,y)$ у кадрі $t$. Результат попарної різниці визначається незалежно для кожного пікселя й каналу:
\begin{equation}
  D_t(x,y,c)=\left|I_t(x,y,c)-I_{t-1}(x,y,c)\right|, \qquad t\geq1.
\end{equation}
Для першого кадру $D_0$ записується чорним, оскільки попереднього кадру немає. Операція має часову складність $O(W\cdot H\cdot C)$ на кадр, але її пікселі не залежать один від одного, тому вона природно відображається на GPU.

CUDA --- платформа NVIDIA для виконання паралельних обчислень на GPU. У цій реалізації OpenCV використано лише для \texttt{VideoCapture} та \texttt{VideoWriter}. Для кожної пари кадрів CUDA Runtime API виділяє буфери пристрою через \texttt{cudaMalloc}, копіює дані через \texttt{cudaMemcpy}, запускає власний kernel і повертає результат у \texttt{cv::Mat} через \texttt{cudaMemcpy}. Декодування й кодування лишаються в OpenCV на CPU, тому вимірюваний час повного конвеєра також містить передавання даних CPU--GPU.

Для систем із кількома GPU застосовано data parallelism за парами кадрів. Вхідний потік формує невеликий пакет послідовних пар кадрів. Для кожного CUDA-пристрою створюється C++ потік, який викликає \texttt{cudaSetDevice()} та обчислює виділені йому пари. Після завершення потоків кадри записуються у вихідне відео в початковому порядку. Такий поділ не порушує залежність між кадрами $I_{t-1}$ та $I_t$, бо обидва кадри передаються одному GPU разом.

\section{Реалізація програми}
Програму реалізовано мовою C++17 у файлі \texttt{src/cuda\_video\_difference.cpp}; власний CUDA kernel міститься у \texttt{src/frame\_difference.cu}. CMake компілює цей файл NVCC і підключає \texttt{CUDA::cudart}; OpenCV потрібний лише з модулями \texttt{core} та \texttt{videoio}.

Алгоритм має такі етапи:
\begin{enumerate}
  \item \texttt{cv::VideoCapture} відкриває вхідний MP4 і зчитує його параметри: ширину, висоту та частоту кадрів.
  \item До результуючого відео записується чорний перший кадр такого самого розміру.
  \item Наступні кадри збираються в пакет, що містить до чотирьох пар на кожен доступний GPU.
  \item Кожен GPU у власному потоці виділяє три CUDA-буфери, копіює туди два кадри та запускає kernel, де один CUDA-потік обробляє один байт кадру.
  \item Результат копіюється в RAM, CUDA-буфери звільняються, а \texttt{cv::VideoWriter} кодує кадри кодеком \texttt{mp4v} у правильному порядку.
\end{enumerate}

Перед початком обробки програма перевіряє число пристроїв викликом \texttt{cudaGetDeviceCount}. За відсутності CUDA-пристрою вона завершується з пояснювальним повідомленням і не видає CPU-результат за GPU-обробку.

\section{Вхідні дані та середовище}
Використано файл \texttt{16756162\_1280\_720\_25fps.mp4} розміром 7\,016\,820 байт, що містить 291 кадр роздільністю 1280$\times$720 із частотою 25 кадрів/с. Локальна машина має NVIDIA GeForce RTX 4070 та драйвер версії 621.01. Для запуску встановлено CUDA Toolkit~13.5 і OpenCV~4.14 з FFmpeg. Звичайний пакет \texttt{opencv-python} із PyPI не застосовується.

\section{Методика запуску та результати}
Програму збирають командами \texttt{cmake -S . -B build} і \texttt{cmake --build build}, після чого запускають із двома аргументами: шляхом до вхідного та вихідного MP4. Перед відкриттям відео виконується перевірка наявності CUDA-пристрою. Якщо OpenCV не має CUDA-підтримки або GPU недоступний, програма завершується з помилкою; реалізація не виконує обчислення різниці на CPU.

Під час одного локального запуску на RTX~4070 створено результуюче відео розміром 3\,409\,817 байт. Виміряний час повного конвеєра становить 2.189~с; він включає запуск програми, декодування, ручні копіювання CPU--GPU, CUDA-обчислення, повернення кадру в RAM і кодування MP4. Пропускна здатність становить $291/2.189=132.937$ кадрів/с. На цій машині доступна одна відеокарта, тому розподіл виконувався одним GPU-потоком; той самий код автоматично розподілить пакети між усіма GPU на багатопристроєвій системі. Ручне виділення та звільнення пам'яті для кожної пари кадрів додає накладні витрати, але демонструє пряме керування CUDA-пам'яттю.

\begin{table}[htbp]
\caption{Результат локального GPU-запуску}
\centering
\begin{tabular}{ll}
\toprule
Параметр & Значення \\
\midrule
GPU & NVIDIA GeForce RTX 4070 \\
CUDA Toolkit & 13.5 \\
OpenCV & 4.14 з FFmpeg \\
Кадрів & 291 \\
Час повного конвеєра & 2.189 с \\
Пропускна здатність & 132.937 FPS \\
Вихідний файл & \texttt{frame\_difference.mp4} \\
\bottomrule
\end{tabular}
\end{table}

Різницевий кадр містить світлі області лише там, де значення пікселів змінилися між двома послідовними кадрами. Статичні ділянки наближаються до чорного, а рухомі об'єкти та зміни освітлення стають помітними.

\KNUUnnumberedSection{Висновки}
Реалізовано програму CUDA/OpenCV для попарної різниці послідовних кадрів відео. Вона відкриває MP4 через \texttt{VideoCapture}, вручну виділяє GPU-пам'ять через \texttt{cudaMalloc}, запускає власний CUDA kernel для різниці кадрів, повертає результат через \texttt{cudaMemcpy} та записує його через \texttt{VideoWriter}. CPU fallback навмисно відсутній: за відсутності CUDA-пристрою програма завершується з помилкою. На RTX~4070 291 кадр оброблено за 2.189~с (132.937 FPS); програма автоматично використовує всі доступні GPU.

\begin{thebibliography}{9}
\bibitem{opencv-cuda} OpenCV Documentation. CUDA-accelerated Computer Vision. \url{https://docs.opencv.org/4.x/d2/dbc/cuda\_intro.html}.
\bibitem{opencv-videoio} OpenCV Documentation. Video I/O module. \url{https://docs.opencv.org/4.x/dd/de7/group\_videoio.html}.
\bibitem{cuda} NVIDIA. CUDA C++ Programming Guide. \url{https://docs.nvidia.com/cuda/cuda-c-programming-guide/}.
\end{thebibliography}

\end{document}
